\documentclass[11pt,letterpaper]{article}
\usepackage[letterpaper,margin=1.0in]{geometry}
\usepackage[T1]{fontenc}
\usepackage{lmodern}
\usepackage{amsmath,amssymb,amsthm}
\usepackage{graphicx}
\usepackage{hyperref}
\usepackage{url}
\usepackage{xcolor}
\usepackage{booktabs}
\usepackage{listings}
\usepackage{enumitem}
\usepackage{microtype}
\usepackage{fancyhdr}
\usepackage{array}
\usepackage{tabularx}
\usepackage{ragged2e}
\usepackage{parskip}

\pagestyle{fancy}
\fancyhf{}
\fancyhead[L]{\small VANTA: An Autonomous Risk Engine for Prediction-Market Credit}
\fancyhead[R]{\small \thepage}
\renewcommand{\headrulewidth}{0pt}

\hypersetup{
  colorlinks=true,
  linkcolor=black,
  citecolor=black,
  urlcolor=blue!45!black,
  pdfborder={0 0 0}
}
\urlstyle{same}

\newtheorem{invariant}{Invariant}

\definecolor{codebg}{rgb}{0.97,0.97,0.97}
\lstset{
  basicstyle=\ttfamily\footnotesize,
  backgroundcolor=\color{codebg},
  breaklines=true,
  frame=single,
  columns=flexible,
  keepspaces=true,
  xleftmargin=4pt,
  xrightmargin=4pt,
  aboveskip=6pt,
  belowskip=6pt,
  showstringspaces=false
}

\setcounter{secnumdepth}{3}

\title{\vspace{-1.0cm}\textbf{\Large VANTA}\\[0.2cm]
\normalsize An Autonomous Risk Engine for Prediction-Market Credit\\
Where Gondor's Curator Sits}

\author{Wisdom O\\[0.12cm]
\small \texttt{okechukwuwisdom016@gmail.com}\\[0.02cm]
\small \texttt{github.com/owizdom}}

\date{April 2026}

\begin{document}

\maketitle
\vspace{-0.4cm}

\begin{abstract}
\noindent
Polymarket crossed \$10B in monthly trading volume and \$500M in open interest by March 2026~\cite{polymarket-volume,polymarket-oi}, but the lending layer Polymarket itself shipped --- PolyLend~\cite{polylend} --- holds \$156. The bottleneck is not appetite; it is risk-management bandwidth. Gondor.fi (\$2.5M pre-seed, December 2025, with participation from Prelude Ventures, Castle Island Ventures, and Maven11~\cite{gondor-seed}) closes part of the gap with a curated beta: 50\% LTV USDC loans against Polymarket ERC-1155 positions, three risk-bucket funds (10\%/20\%/30\% interest caps), markets hand-selected and rebalanced by the team.

VANTA is the autonomous risk engine that sits where Gondor's curator sits. Same product, different brain. The brain is a sovereign agent running three continuous reasoning loops --- credit, model, operational (\S\ref{sec:loops}) --- above a deterministic floor of formula-priced haircuts (\S\ref{sec:haircut}) and contract-enforced onboarding gates (\S\ref{sec:onboard}). The floor sets what the contract will accept; the reasoning sets what the agent proposes within that envelope --- per-market caps inside the \$25k probation ceiling, calibration-drift flags to lender governance, per-loan health monitoring, cluster-correlation alerts above the explicit tag system. Every reasoning step is a typed signed event with full input provenance; the chain-of-thought is auditable and bounded, not opaque. The agent runs inside an EigenCompute Intel TDX enclave; its Ed25519 signing key is generated in enclave RAM at boot and never leaves. The on-chain contracts (\texttt{LpVault}, \texttt{LoanBook}, \texttt{VantaVault}) are owned by an address only the running enclave image can derive. The team's role shifts from ongoing fund manager to one-time image publisher; trust shrinks from operational to one-time auditor review.

We make four claims, in descending strength of evidence. First, the addressable market splits cleanly: a \emph{formalizable subset} of Polymarket (standard templates, whitelisted resolvers, $\geq 14$ days old) where curation reduces to gates plus bounded reasoning, and a \emph{hard subset} (novel templates, correlated baskets, custom resolvers) where human curation remains the correct tool. VANTA owns the formalizable subset; we explicitly defer the hard subset to lender-governance upgrades, not to the agent. Second, the resolution-aware haircut is a more honest pricing of tail risk than three buckets on the markets where its inputs are well-defined. Third, continuous agent reasoning --- not a periodic script over hard-coded gates --- is the correct shape for live risk: collateral prices, dispute behavior, and cluster correlations evolve between gate cycles, and the three loops respond to that evolution above the contract floor. Fourth, an attestation-liveness exit gives lenders a custody guarantee Gondor cannot offer: withdraw-only mode that activates without admin signatures if the enclave goes silent.

This paper specifies the engine, the onboarding rule set, the threat model, and the boundaries of the autonomy claim --- including the places where the framing overreaches and what we substitute instead.
\end{abstract}

\vspace{0.3cm}

% =========================================================
\section{Where VANTA Sits}
\label{sec:ladder}

The Polymarket lending stack today has three live layers, in increasing order of risk-management sophistication.

\textbf{PolyLend} is Polymarket's own peer-to-peer lending product, deployed from the Polymarket repository~\cite{polylend}. As of April 2026 it holds \$156 in TVL on DefiLlama. The product Polymarket shipped is dormant. The reason is not infrastructure: the contracts work. The reason is risk management. Peer-to-peer means each lender prices their own loan, and few lenders have the appetite to compute the haircut math on a 2026-midterms YES position alone.

\textbf{Gondor.fi} closes that gap with a curated beta. Borrow up to 50\% LTV against Polymarket ERC-1155 positions deposited into Morpho-deployed vaults~\cite{morpho}; three lender-side funds (Conservative, Moderate, Growth) with interest caps of 10\%, 20\%, and 30\%; and --- the core --- markets hand-selected by the team based on liquidity, spreads, duration, resolution criteria, and slippage~\cite{gondor-docs}. Funds are actively managed: the team adds new supported markets and rebalances liquidity on live conditions. Risk decisions live in the team, not the contracts. The closest thing today, and the closest thing to product-market fit. The ceiling is curator bandwidth.

\textbf{VANTA} is what Gondor would be if the curator were inside the contract. The product is the same: USDC loans against Polymarket positions, with a haircut, with liquidations, with an LP vault. The brain is different. Every loan is priced by a continuous resolution-aware formula (\S\ref{sec:haircut}). New markets are onboarded by an agent running inside an enclave (\S\ref{sec:onboard}, \S\ref{sec:loops}), gated by hard contract limits and bounded by a 3-markets-per-day throughput cap and a per-market \$25k probation envelope. Parameter changes are signed events written by the same enclave (\S\ref{sec:verify}). The team publishes the image; the auditor checks the image; the running code is what the image was.

\textbf{What we are not claiming.} It is tempting to read VANTA as ``Gondor automated.'' That overstates and misframes. Gondor's curation is not only bandwidth-limited; on its hardest markets, curation \emph{is} the product --- the team's accumulated knowledge of which resolvers misbehave under stress, which sub-markets are one cluster bet, which custom resolution texts hide ambiguity. A formula and an LLM gate cannot substitute for ten years of watching UMA disputes resolve. We do not pretend they can.

What we claim instead is a clean split. Polymarket has a \emph{formalizable subset} --- standard-template markets on whitelisted resolvers above a 14-day age floor --- where the gates of \S\ref{sec:onboard} are well-defined and the reasoning loops of \S\ref{sec:loops} have legible inputs. On that subset, the curator-bandwidth ceiling is the binding constraint: a team of three can curate fifteen markets well; an agent reasoning under contract-enforced caps can clear fifteen hundred. The remaining \emph{hard subset} --- novel templates, custom resolvers, deeply correlated baskets --- stays in human hands by lender-governance upgrade, not by agent fiat. VANTA does not eat the hard markets. It removes the bandwidth ceiling from the markets where the math has well-defined inputs, and leaves the rest to the curators who already do that work better than any model will. bobIsAlive~\cite{bobisalive} and Sovra~\cite{sovra} are the substrate (sovereign agents that hold their own keys and pay their own bills); VANTA is the use case (a credit protocol where the agent's reasoning is bounded, signed, and verifiable above a contract floor).

% =========================================================
\section{The Problem: Locked Capital, Curator-Bound Credit}
\label{sec:problem}

Polymarket positions are time-locked claims on USDC. A YES share at \$0.40 on a market resolving in 90 days is a 90-day lockup of \$0.40 of capital with a Bernoulli payoff. For traders sized like Th\'eo --- the 2024 election whale who Chainalysis~\cite{theowhale} reconstructed at \$85M across 11 accounts, sitting \$3M underwater the day before the election --- ``locked'' is not a metaphor. The position is right; the capital cannot be redeployed.

Th\'eo is the upper edge. The borrower VANTA is designed for is the mid-tier directional holder: a \$10k--\$5M position on a 2026-midterms YES, a 2028-presidential bundle, an FOMC rate-path basket, or a 6--18-month geopolitics market. They could exit by selling into Polymarket's CLOB at a discount; they could close by buying the opposite outcome; they prefer to do neither, because the position is correct on their model and the alternative use of unlocked capital --- another Polymarket position, a stablecoin yield, a mainnet trade --- has an opportunity cost the discount does not cover. Today they hold to expiry (capital locked), sell into a thin book on long-tenor markets (price the discount), or queue for Gondor's hand-curated whitelist (waitlist, not all markets). PolyLend's \$156 TVL~\cite{polylend-tvl} is the revealed preference: the peer-to-peer alternative is not solving for them. The bound on this persona is the time horizon: VANTA's 14-day market-age floor and the $\tau_{\mathrm{gap}}$ term in the haircut formula (\S\ref{sec:haircut}) are sized for medium-to-long-tenor positions, not for the short-resolution speculator who can afford to wait out the lockup.

The aggregate scale is no longer marginal. Polymarket reported \$10.57B in March 2026 monthly volume~\cite{polymarket-volume}, the first time crossing \$10B. Open interest crossed \$500M by April 2026~\cite{polymarket-oi}, having grown from roughly \$212M in October 2025. Sector-wide open interest --- Polymarket plus Kalshi plus the long tail --- crossed \$1B in the same window~\cite{predictionmarkets-1b}. Long-horizon markets (2026 midterms, fed-rate trajectories, geopolitics) tie capital up for weeks or months at a time.

The lending side has not kept pace. PolyLend is dormant at \$156 TVL~\cite{polylend-tvl}. Gondor is in beta and supplied lender liquidity from the team's own treasury. Kalshi's affiliate Kinetic Markets received NFA approval as a Futures Commission Merchant in March 2026~\cite{kalshi-fcm}, clearing the path for margin trading on event contracts under CFTC regulation; that is real competition, but it does not unlock pre-resolution capital from existing Polymarket positions, and it does not solve the problem for the on-chain trader.

Two specific frictions explain the gap, both load-bearing for VANTA's design:

\begin{enumerate}[leftmargin=*,itemsep=3pt,topsep=3pt]
  \item \textbf{Per-market risk math.} Each Polymarket market has its own resolution date, dispute history, depth profile, and resolver. A flat LTV ratio over-reserves on liquid near-resolution markets and under-reserves on illiquid far-from-resolution ones. The math has to be per-market or it is wrong.
  \item \textbf{Per-market onboarding cost.} Reviewing a new market for liquidity, spread, duration, resolution criteria, and slippage is hours of human work. A team curating fifteen good markets cannot curate fifteen hundred. Polymarket has thousands of active markets at any given moment.
\end{enumerate}

A formula that does the per-market math, executed by a process that scales without human review, and gated such that ``scales without review'' does not mean ``onboards garbage,'' is the sufficient condition. The next sections specify each piece.

% =========================================================
\section{Invariants}
\label{sec:invariants}

These properties must hold for VANTA to be what it claims.

\begin{invariant}[Agent-owned identity]
VANTA's cryptographic identity is an Ed25519 keypair generated in enclave RAM at boot. The public key is the agent's canonical name. The private key never leaves the enclave. No founder, cloud admin, or steward holds this key.
\end{invariant}

\begin{invariant}[Signing separation]
The enclave signs a risk-model attestation binding the loan's terms to the model that produced them. The borrower's wallet signs the loan acceptance. The vault contract validates both before releasing USDC. No single key originates a loan.
\end{invariant}

\begin{invariant}[Public provenance]
Every credit event, model event, onboarding event, and operational event is a typed JSON record signed by the enclave. The event log is content-addressed, pinned to Walrus and IPFS, and tip-anchored weekly to Base.
\end{invariant}

\begin{invariant}[Bounded autonomy envelope]
The agent's onboarding authority is bounded by hard limits the contract enforces, not by self-imposed budgets: $\leq 3$ new markets per 24h, $\leq \$25\text{k}$ probationary borrow capacity per new market for the first 7 days, UMA-resolver whitelist, resolution-text-template whitelist (\S\ref{sec:onboard}). New limits require a 7-day timelocked lender-quorum upgrade.
\end{invariant}

\begin{invariant}[Attestation-liveness exit]
If the contract receives no fresh attestation heartbeat for 14 days, it enters Withdraw-Only mode: no new loans, no parameter updates, no onboarding. Depositors withdraw pro-rata against on-chain accounting state at TWAP-marked NAV. The TEE admin key is irrelevant in Withdraw-Only mode --- withdrawals are permissionless. Custody never depends on a single live infrastructure.
\end{invariant}

\begin{invariant}[Self-sustaining economics]
Operating costs --- enclave hosting, LLM inference, gas, oracle reads --- are paid from the protocol's own treasury, funded by loan-origination and oracle fees. Inflows and outflows are signed events; the treasury balance is readable on-chain.
\end{invariant}

% =========================================================
\section{Pricing: Resolution-Aware Haircuts}
\label{sec:haircut}

A binary prediction-market price is bounded in $[0,1]$, has a terminal date, and is dominated by time-to-resolution volatility. A flat LTV ratio over-reserves far-from-resolution and under-reserves near-resolution. Three discrete buckets discretize the same risk surface; they are legible and computationally cheap, but they leave money on the table on liquid mid-tenor markets and bear too much risk on long-tail-text markets near resolution.

We price the collateral value of a YES position as follows. Let $p$ be the market price, $\tau$ the time to resolution, and $\rho$ the oracle dispute rate of similar markets on the same resolver:

\begin{equation}
\label{eq:haircut}
h(p, \tau, \rho) = \alpha \, \sigma(\tau) \sqrt{\tau} \;+\; \beta (1-p) \, e^{-\tau / \tau_{\mathrm{gap}}} \;+\; \gamma \, \rho
\end{equation}

The applied haircut is clamped: $h_{\mathrm{applied}} = \min(h, 1 - \varepsilon)$, $\varepsilon = 0.05$. Collateral value $V = p \cdot (1 - h_{\mathrm{applied}}) \cdot N$ for notional $N$. The implied LTV ceiling is bounded above by Gondor's 50\% --- VANTA is never more aggressive than Gondor on the LTV axis --- and below by the safe floor implied by the formula.

Each term answers a distinct question. \textbf{Drift} $\alpha \sigma(\tau) \sqrt{\tau}$ asks how much the price might move before resolution; $\sigma(\tau)$ is calibrated per the jump-diffusion framework of Kirilenko et al.~\cite{kirilenko2025}, not $\sqrt{p(1-p)}$ (which double-counts terminal risk). \textbf{Gap} $\beta(1-p)e^{-\tau/\tau_{\mathrm{gap}}}$ asks how much room there is for the position to fall to zero, weighted heavily as resolution approaches. \textbf{Oracle} $\gamma \rho$ asks what fraction of similar markets get disputed; UMA's published dispute rate is in the 1.3--1.67\% range~\cite{uma-disputes}, with a whitelisted-proposer regime introduced in August 2025 covering 96\% of Polymarket proposals at 99.7\% accuracy~\cite{uma-managed}.

Baselines: $\alpha = 1.0$, $\beta = 0.6$, $\gamma = 1.0$, $\tau_{\mathrm{gap}} = 7$ days, $\varepsilon = 0.05$. We pin $\sigma(\tau)$ to two anchors: $\sigma\sqrt{\tau} = 0.15$ at $\tau = 365$ days, $\sigma\sqrt{\tau} = 0.04$ at $\tau = 14$ days. The closed-form fit is $\sigma(\tau) = 0.013724 \cdot \tau^{-0.094658}$, clamped to $[0.005, 0.05]$.

\textbf{Worked example.} YES on a 2028 election market at $p=0.60$, $\tau=365$, $\rho=0.02$:
\begin{align*}
\text{drift} &= 1.0 \cdot 0.013724 \cdot 365^{-0.0947} \cdot \sqrt{365} \approx 0.150 \\
\text{gap}   &= 0.6 \cdot 0.40 \cdot e^{-365/7} \approx 0 \\
\text{oracle}&= 1.0 \cdot 0.02 = 0.020 \\
h            &\approx 0.170
\end{align*}
On a 250{,}000 USDC notional, collateral value is $0.60 \cdot 0.83 \cdot 250{,}000 = 124{,}500$ USDC. A 60{,}000 USDC loan against it is 48.2\% LTV. Under Gondor's Growth fund (30\% interest cap), the same notional would qualify only as a 50\% LTV bucket if the team chose to onboard it; the formula tells you what the right haircut actually is.

\textbf{What the formula does not do.} It does not read resolution-criteria text for ambiguity, model correlated markets, or catch a flash glitch in the CLOB feed. Each of these is handled by a separate gate (\S\ref{sec:onboard}) or circuit breaker (\S\ref{sec:threats}). The formula is the pricing layer; the gates are the eligibility layer. Combining them is what delivers ``scales without curating garbage.''

% =========================================================
\section{Origination: Two Signatures, Post-Image}
\label{sec:settlement}

Originating a loan requires two signatures. The borrower signs an EIP-712 acceptance binding $(\text{loanId}, \text{principal}, \text{haircutBps}, \text{maturityTs})$. The runtime, holding the agent's HKDF-derived secp256k1 admin key, broadcasts \texttt{LoanBook.originate(...)} on Base Sepolia. The contract recomputes and stamps $\text{paramsHash} = \text{keccak256}(\text{abi.encode}(\text{borrower}, \text{principal}, \text{haircutBps}, \text{maturityTs}))$ on the loan record.

Only \emph{after} the broadcast confirms at depth $\geq 2$ does the runtime sign the \texttt{loan.origination} event into the log, with \texttt{(txHash, blockNumber, blockHash, paramsHash)} embedded. This is the post-image pattern: a crash mid-flight leaves no phantom origination, because the event isn't written until the chain has confirmed.

The verifier (CLI: \texttt{vanta-verify <url>}) cross-checks every \texttt{loan.origination} event against on-chain \texttt{LoanBook.loans(loanId).paramsHash}. Mismatch is a hard fail. Same guarantee whether the verifier runs on the lender's laptop or in CI.

\textbf{Pledge mechanics.} VANTA supports Polymarket as its first venue. Polymarket users hold CTF positions through Gnosis Safe proxy wallets~\cite{polymarket-proxy}. Primary path: borrower's proxy calls \texttt{safeTransferFrom} into VantaVault on Polygon Amoy for the loan duration --- the PolyLend pattern~\cite{polylend}. Secondary path: revocable approval (\texttt{setApprovalForAll(VantaVault, true)}); priced at $5\times$ the escrow haircut floor.

\textbf{Settlement.} On the underlying market's resolution, the pledged position pays out via Polymarket's CTF; \texttt{LoanBook.markRepaid} pulls principal back into LpVault via OpenZeppelin's \texttt{SafeERC20.safeTransferFrom}; the residual is the borrower's. \texttt{markLiquidated(loanId, auctionEscrow, proceeds)} is the analogous path for forced settlement.

\textbf{Liquidation mechanism.} VANTA inherits the YES--NO merge pattern from Gondor verbatim: at the 77\% LTV liquidation threshold, seize 77\% of collateral, simultaneously buy opposite-outcome shares on Polymarket, merge YES--NO pairs, redeem at \$1. Polymarket's order books are off-chain, so hedging via the opposite side minimizes liquidation latency relative to selling collateral into a thin book. We do not improve on this; we copy it. The honest contribution here is execution-time verifiability, not mechanism novelty.

\textbf{Borrower flow, single-position v1.} The minimal v1 path is single-loan: the borrower's proxy pledges a position via \texttt{safeTransferFrom} into VantaVault, signs the EIP-712 acceptance, receives USDC, and either repays before resolution to reclaim the position or lets settlement clear automatically. The secondary revocable-approval path (\texttt{setApprovalForAll(VantaVault, true)}, priced at $5\times$ the escrow haircut floor) supports re-borrow against the same already-approved position without a second \texttt{safeTransferFrom}, and lets a borrower draw against a portfolio of approved positions in sequence.

\textbf{Borrower flow, multi-position v1.x.} A holder of ten Polymarket positions in v1 must originate ten loans separately to monetize all ten; there is no batched origination, no portfolio dashboard, and no one-click ``leverage up an existing pledge'' on top of an already-active loan. Exit on resolution is automatic per market but does not collapse a multi-loan portfolio into a single settlement transaction. These are real frictions that adjacent designs --- USDC-funded order-book leverage with conservative LLTVs, surveyed in lender conversations --- avoid by collapsing pledge and origination into a single user action. We treat multi-position UX as v1.x scope (\S\ref{sec:roadmap}, Month 5--6). VantaVault already tracks position-level escrow, so the contract surface is in place; the missing pieces are a batched-origination event schema, a position-portfolio view served from the same signed event log the verifier reads, and a credit-loop-aware exit planner (\S\ref{sec:loops}) that schedules correlated settlements together rather than serially. The reasoning-loop architecture matters here: the credit loop already monitors per-loan health continuously, so a portfolio-level exit planner is a query against state the agent already maintains, not a new subsystem.

% =========================================================
\section{Autonomous Market Onboarding}
\label{sec:onboard}

This is the load-bearing differentiation against Gondor. If onboarding has to be human-curated, VANTA is Gondor with a fancy admin key. If onboarding is autonomous and bounded, VANTA is a different protocol.

\textbf{Two layers.} Onboarding has a contract-enforced \emph{gate floor} that defines eligibility, and bounded \emph{agent reasoning} above the floor that sets per-market parameters within the autonomy envelope. The floor is what the contract will ever accept. The reasoning is what the agent contributes between cycles. This section specifies both.

The agent runs an onboarding cycle every six hours. For each candidate market on Polymarket, the runtime collects seven inputs:

\begin{itemize}[leftmargin=*,itemsep=2pt,topsep=2pt]
  \item \texttt{depth\_5pct}: USDC notional needed to move best-bid by 5\% (CLOB snapshot).
  \item \texttt{dispute\_30d}: count of UMA disputes against this market's resolver in the last 30 days.
  \item \texttt{age\_days}: days since market creation.
  \item \texttt{vol\_7d}: 7-day realized volatility of mid-price.
  \item \texttt{text\_hash}: keccak256 of the full resolution-criteria text.
  \item \texttt{text\_score}: deterministic LLM-judge output (0--100) on resolution-criteria \emph{clarity}, with the prompt fixed and committed in the enclave image.
  \item \texttt{tags[]}: pre-registered correlation tags (e.g.\ \texttt{us-elections-2028}, \texttt{fomc-2026}).
\end{itemize}

Every gate must pass for the market to be onboarded:

\begin{table}[h]
\centering
\small
\caption{Onboarding gates and thresholds. All must pass.}
\label{tab:gates}
\begin{tabularx}{\textwidth}{@{}p{0.22\textwidth}p{0.18\textwidth}X@{}}
\toprule
\textbf{Gate} & \textbf{Threshold} & \textbf{Why} \\
\midrule
Depth & $\geq \$250\text{k}$ at 5\% impact & Liquidation can clear without crossing more than the buffer in \S\ref{sec:settlement}. \\
Disputes (30d) & $0$ on resolver & A disputed resolver in the last 30 days is a 30-day waiting room. \\
Age & $\geq 14$ days & New markets attract toxic flow before depth stabilizes. \\
7-day volatility & $\leq 0.35$ price units & Above this, the formula's $\sigma(\tau)$ pin is unreliable. \\
Text score & $\geq 70$ & Resolution-criteria clarity below this is a flag, not a feature. \\
Text template & Hamming distance $\leq 32$ & Match against a fixed library of Polymarket standard templates; custom-text markets default-reject. \\
Tag novelty & $\subseteq$ pre-registered set & New tag families require a separate timelocked governance upgrade, not an onboarding. \\
\bottomrule
\end{tabularx}
\end{table}

\textbf{Reasoning above the floor.} Passing the seven gates of Table~\ref{tab:gates} is necessary, not sufficient. The agent then reasons --- inside the enclave, with the reasoning trace itself a signed event --- over the inputs in their \emph{full} form, not just the digests the contract checks: the resolution-criteria text against historical resolver behavior on similar phrasings; the depth shape over the trailing 30 days, not just the snapshot; the timing of recent UMA disputes on this resolver against the candidate's own resolution date; cluster-correlation evidence across the agent's existing exposure; a deterministically-prompted LLM judge whose output is bounded to a structured score with the prompt fixed and committed in the enclave image. The reasoning produces three outputs the contract reads: \texttt{decision} (onboard, reject, or flag for human review), \texttt{cap\_initial\_usdc} within the \$25k probation ceiling, and \texttt{ltv\_max\_bps} within the formula's safe-floor-to-50\% band. The full reasoning trace --- inputs, intermediate scores, model identifier, prompt hash, decision rationale, dissenting considerations --- is committed in a sibling \texttt{ReasoningTrace} event signed by the same enclave key and pinned alongside the \texttt{OnboardEvent}. The contract does not read the trace; lenders, auditors, and the verifier CLI do.

\textbf{Why reasoning rather than a fixed rule.} A purely rule-based onboarder fails on the part of a Polymarket market that is not numeric: the resolution-criteria text. ``Trump wins the 2028 election'' and ``Trump is inaugurated as the 47th President on January 20, 2029'' resolve to the same nominal outcome but admit different oracle-precision risks --- one is a calendar date the resolver checks against AP coverage, the other depends on whether ``wins'' means popular vote, electoral college, or a network call. Hamming distance against template families catches wholesale custom text but not these template-fitting-but-semantically-different cases. The \texttt{text\_score} gate catches gross unclarity; the agent's reasoning catches the long tail. The same logic applies to depth-shape evolution between snapshots, dispute-timing patterns the 30-day count flattens, and cluster correlations the pre-registered tag set cannot anticipate. Where math stops and reasoning starts is precisely here: the formula and the gates handle inputs whose shape is closed; the agent handles the inputs whose shape is not. Both layers are bounded above by the same hard contract limits.

\textbf{Safety floor.} Three layers of bound make ``scales without curating garbage'' actually true rather than aspirational:

\begin{enumerate}[leftmargin=*,itemsep=2pt,topsep=2pt]
  \item \textbf{Throughput cap.} The contract rejects more than 3 \texttt{OnboardEvent}s per 24h, regardless of how many candidates pass gates. Even a maximally adversarial gate-pass storm cannot onboard a fourth.
  \item \textbf{Blast-radius cap.} Each newly-onboarded market enters a 7-day probation period during which the per-market borrow capacity is hard-capped at \$25k. The cap ramps to formula-derived after probation if no flags or disputes triggered. The maximum harm a single onboarding decision can do is \$25k.
  \item \textbf{Whitelist envelope.} Resolver address must be in a hardcoded UMA-resolver whitelist; resolution text must match a pre-registered Polymarket template family. Adding new resolvers or template families is a separate governance action --- timelocked, lender-quorum, off the agent's authority.
\end{enumerate}

\textbf{Signed event schema.} Each onboarding decision is a signed event the contract validates:

\begin{lstlisting}
OnboardEvent {
  market_id:        bytes32,
  attestation:      bytes,        // EigenCompute TDX quote, fresh
  inputs_hash:      bytes32,      // keccak256(depth || disputes || age || vol
                                  //           || text_hash || tags)
  trace_hash:       bytes32,      // keccak256(ReasoningTrace event payload)
  decision:         bool,
  ltv_max_bps:      uint16,
  rate_params:      (alpha, beta, gamma, tau_gap),
  cap_initial_usdc: uint128,      // <= 25_000 during probation
  timestamp:        uint64,
  signature:        bytes65       // ed25519 from in-enclave key
}
\end{lstlisting}

The contract honors \texttt{OnboardEvent} only if the signature verifies against the address derived from the published image's measurement, the attestation is fresh, and \texttt{block.timestamp - timestamp < 5 min}. Stale or malformed events revert.

\textbf{Fail-safes.} Bounded autonomy is not a substitute for emergency response. Three escalation paths exist, in increasing severity:

\begin{itemize}[leftmargin=*,itemsep=2pt,topsep=2pt]
  \item \textbf{Per-market freeze.} Any depositor posts a \$1k bond against a market; if 3+ independent flags arrive within 24h, the market enters frozen state --- no new borrows, existing positions repay or liquidate normally. Bond returned if the flag is upheld by a 7-day post-mortem window; slashed if frivolous.
  \item \textbf{Global onboarding pause.} A 4-of-7 lender-quorum multisig can halt new \texttt{OnboardEvent}s for 7 days. The multisig cannot move funds, only pause onboarding. This preserves ``no humans in custody'' while permitting ``humans on the emergency brake.''
  \item \textbf{EigenLayer slashing.} The EigenCompute operator stake backing the enclave is slashable if the attestation chain breaks or if the published image disagrees with the running measurement. Provides real economic backing for the verifiability claim, not just cryptographic.
\end{itemize}

This is bounded autonomy, not unbounded. Calling the design ``fully autonomous'' overstates it; the more honest framing is \emph{verifiable autonomy within a lender-governed envelope}, where the envelope itself is timelocked and on-chain.

% =========================================================
\section{Continuous Reasoning Loops}
\label{sec:loops}

The onboarding cycle (\S\ref{sec:onboard}) is one of three continuous loops the agent runs inside the enclave. Onboarding is the loop the contract gates most strictly because its blast radius is the largest --- a bad onboarding decision shapes exposure for weeks. The other two --- credit and operational --- run continuously between onboarding cycles and carry the bulk of the agent's cognitive load over a typical month. This is the difference between VANTA and a periodic script: between gate cycles, the protocol is not idle.

\textbf{Credit loop.} The agent monitors every active loan against the live state of its collateral. Inputs: the current best-bid and 30-minute TWAP from the Polymarket CLOB, the market's UMA dispute history since origination, the depth at the 5\% impact threshold, the time remaining to resolution. The job is to flag positions whose collateral value is approaching the 77\% liquidation threshold before the contract's hourly cap-recompute does, so that the borrower can be notified, voluntary repayment can clear without auction, and the auction-spread floor (\S\ref{sec:threats}, T6) does not become the realized exit. The loop also reasons about cluster events: when a 2028-election sub-market receives a resolution news shock (e.g., a candidate withdrawal), every loan in the cluster is repriced and a per-market freeze can be requested before the depositor-bond mechanism arrives at the same conclusion through three independent flags.

\textbf{Model loop.} The haircut formula (\S\ref{sec:haircut}) has a closed form, but its parameters $(\alpha, \beta, \gamma, \tau_{\mathrm{gap}})$ are calibrated against a public dataset of resolved markets. The model loop runs adversarial replay weekly: synthetic borrowers attempt to select into mispriced loans across a generated stress portfolio (election shocks, dispute clusters, oracle-resolution disagreements), and the agent records realized loss-given-default versus the formula's pre-loss prediction. When realized error exceeds a calibrated band, the loop emits a \texttt{CalibrationProposal} event carrying the proposed parameters, the dataset slice that triggered it, and the realized-vs-predicted distribution. The proposal is \emph{not} self-actuating: it is an input to the 4-of-7 lender-quorum multisig, gated by the 7-day timelock of T8. The agent reasons; the lenders vote; the contract delays.

\textbf{Operational loop.} The agent pays its own bills (Invariant~6). The treasury wallet (\texttt{vanta-treasury-v1}) receives oracle and origination fees; the operational loop reasons over runway, gas-price evolution on Base and Polygon Amoy, EigenCompute hosting renewals, LLM inference cost, and oracle-read amortization. When projected runway falls below a defined threshold, the loop emits a \texttt{TreasuryAlert}; when costs spike anomalously (gas above the 95th percentile of the trailing 30 days, inference cost step-change), the loop emits an \texttt{OperationalAnomaly}. Neither event triggers a contract action; both are inputs to the lender multisig and to the verifier audit. The agent's own economic survival is part of its reasoning surface.

\textbf{What the loops are not.} They are not autonomous fund management. None of the three loops can move funds, change formula parameters in production, or onboard markets outside the \S\ref{sec:onboard} envelope. The credit loop can request a per-market freeze; only the depositor-bond mechanism executes one. The model loop can propose recalibration; only the multisig's timelocked quorum can apply one. The operational loop can flag treasury risk; only governance can respond. The reasoning is bounded at every layer.

\textbf{Reasoning provenance.} Every loop emits typed JSON events signed by the in-enclave Ed25519 key (Invariant~3). The schema for a credit-loop iteration:

\begin{lstlisting}
CreditLoopTick {
  loan_id:           bytes32,
  collateral_value:  uint256,    // V = p * (1 - h_applied) * N
  ltv_current_bps:   uint16,
  trace_hash:        bytes32,    // sibling ReasoningTrace event
  flag:              uint8,      // 0 = ok, 1 = watch, 2 = freeze-request
  timestamp:         uint64,
  signature:         bytes65
}
\end{lstlisting}

\texttt{vanta-verify} walks every loop's event stream alongside the on-chain origination and settlement events, cross-checking the agent's reasoning trace against the actions the contract observed. A loan that liquidated without a prior credit-loop watch flag is auditable; an onboarding that occurred without a corresponding reasoning trace is rejected by the verifier as malformed.

% =========================================================
\section{Verifiable Execution}
\label{sec:verify}

VANTA runs inside one EigenCompute Intel TDX enclave. At boot, the runtime calls the \texttt{@layr-labs/ecloud-sdk} \texttt{AttestClient} to exchange a hardware quote for an encrypted JWT bound to the enclave's measured identity. The enclave receives a \texttt{MNEMONIC} secret through this attested channel. From that one seed, HKDF-SHA-256 derives:

\begin{itemize}[leftmargin=*,itemsep=2pt,topsep=2pt]
  \item \texttt{vanta-signing-v1} --- the Ed25519 signing key used on every event.
  \item \texttt{vanta-treasury-v1} --- secp256k1 wallet on Base that pays the agent's bills.
  \item \texttt{vanta-origination-eoa} --- secp256k1 admin EOA on \texttt{LpVault} and \texttt{LoanBook}.
  \item \texttt{vanta-vault-operator-v1} --- lender-vault operator signature.
  \item \texttt{vanta-api-credentials-v1} --- encryption key for in-enclave API credentials.
\end{itemize}

All five derivations happen inside the enclave; none of the derived keys leave it. The current deployments are on testnets: \texttt{LpVault} at \texttt{0xD9AFC3abC151239B047F653e0493792ecEB0f076} on Base Sepolia, \texttt{LoanBook} at \texttt{0xA6004cb25ce50baEDa38017B81577FA2419a876f} on Base Sepolia, \texttt{VantaVault} at \texttt{0xD9AFC3abC151239B047F653e0493792ecEB0f076} on Polygon Amoy. All three are owned by \texttt{0x8275001D7618263304B76Ed621150cE4193826e6}, the address derived inside the enclave. EigenCompute App ID is \texttt{0xaa4875Ff1514b8227Dc61466473c0B1E4C83CAf9}.

\textbf{What ``verifiable'' actually means here.} Eigen's attestation chain proves the running code's measurement matches the published image's measurement, byte-for-byte. The depositor checks this in one of two ways. The honest way: an auditor publishes a build of the image, derives the measurement, and signs a statement that the measurement equals the running attestation. The depositor trusts the auditor, not the team. The aspirational way: the depositor reproduces the build themselves. We do not pretend many depositors will. The trust shift VANTA actually provides is from \emph{ongoing operational trust} (Gondor: ``the team uses the feed honestly today and tomorrow and next month'') to \emph{one-time auditor review} (VANTA: ``the auditor checked the image; the running code is what the image was''). That is a real reduction; it is not the elimination of trusted parties.

\textbf{Sovereignty gap.} EigenCompute's KMS today is a single-vendor (GCP Confidential Space) HMAC root: whoever holds \texttt{upgradeApp} can mint a new image with the same HKDF-derived secrets. Mitigation: the Ed25519 signing key is freshly generated per boot and is \emph{not} HKDF-derived from \texttt{MNEMONIC}, so a malicious image fork breaks the signature chain visibly in the event log. The attestation-liveness exit (Invariant~5) makes this a recoverable failure mode, not a stranded-funds one. Stewards multisig + 7-day timelock on \texttt{upgradeApp}; conservative TVL ceiling until single-vendor KMS is replaced.

\textbf{What the verifier actually checks.} The off-chain CLI \texttt{vanta-verify} reads \texttt{/api/events/:id}, walks the chain to genesis, verifies each Ed25519 signature with strict RFC 8032, and for every \texttt{loan.origination} or \texttt{loan.onboarding} event cross-checks the on-chain \texttt{paramsHash}. Eight artifacts a thorough verifier checks: the agent's signing public key, the EigenCompute attestation JWT, the current live risk model and history, the signed event log, the agent's treasury state, the on-chain \texttt{LpVault} state, the per-market onboarding cap registry, and the published image build instructions.

% =========================================================
\section{Comparison: Gondor's Curator vs. VANTA's Agent}
\label{sec:compare}

Where the contrast actually lands, and where it does not.

\begin{table}[h]
\centering
\footnotesize
\caption{Mechanism-by-mechanism comparison. Honest assessment of where VANTA's autonomous design wins, ties, or trades.}
\label{tab:compare}
\begin{tabularx}{\textwidth}{@{}p{0.18\textwidth}p{0.26\textwidth}X@{}}
\toprule
\textbf{Mechanism} & \textbf{Gondor (manual)} & \textbf{VANTA (autonomous)} \\
\midrule
Market onboarding & Team hand-selects per liquidity, spread, duration, resolution criteria, slippage. & 6h cycle, 7 inputs, 7 hard gates as floor, bounded agent reasoning above floor (\S\ref{sec:loops}), 3-markets/day cap, \$25k probation envelope. Caps at enclave bandwidth on the formalizable subset. \emph{Trade}: hard subset (custom text, novel resolvers, deeply correlated baskets) is explicitly deferred to lender-governance upgrade, not eaten by the agent. \\
Interest-rate caps & Three buckets: 10\%, 20\%, 30\%, fixed per market. & Continuous formula; rate is a function of the same inputs that drive the haircut. \emph{Trade}: legibility --- buckets are easy to reason about; formula is not. Mitigation: publish realized-rate histograms per market. \\
LTV ceiling & Fixed at 50\%. & Per-market ceiling derived from formula, capped at 50\% (never above), floored at the safe floor. Variable but bounded. \\
Liquidation threshold & 77\%. & 77\%, identical. No edge. \\
Liquidation mechanism & Seize, hedge via opposite outcome, merge YES--NO, redeem at \$1. & Identical mechanism. Polymarket-structural, not curator-vs-agent. \\
Collateral pricing & $\min(\text{best-bid}, \text{TWAP})$, computed by team pipeline. & $\min(\text{best-bid}, \text{TWAP})$, computed by attested code. Same formula, verifiable execution. \\
Exposure caps & Per-market dollar caps, set by team. & Per-market caps recomputed hourly from realized depth $\times k$; cluster correlations monitored continuously by the credit loop (\S\ref{sec:loops}) above the tag-system floor. \emph{Trade}: novel correlations covered by neither layer remain weaker than a domain-expert curator (\S\ref{sec:threats} T7). \\
Early closure & Linear 77\%$\to$0\% LTV ramp over the last 7 days. & Identical ramp; haircut formula's $\tau_{\mathrm{gap}}$ anticipates the ramp by raising rates earlier. \\
Fund segmentation & Three funds (Conservative/Moderate/Growth). & Single pool. \emph{Trade}: no off-the-shelf risk-segregated SKU for retail. Mitigation: senior/junior tranches as future work. \\
Cross-margining & V1 (2026) with curator-defined bundles. & Bundle eligibility from the same formula with correlation haircut. Parity at best until both ship. \\
Custody & Funds in Morpho-deployed vaults; admin keys held by team multisig. Recoverable. & Funds in TEE-owned contracts; admin = enclave-derived address. \emph{Trade}: not multisig-recoverable. Mitigation: attestation-liveness exit (Invariant~5) gives permissionless withdraw-only after 14 days of silence. \\
\bottomrule
\end{tabularx}
\end{table}

The contest is not a sweep. VANTA cleanly improves onboarding throughput, exposure-cap responsiveness, pricing granularity, and custody verifiability. VANTA matches Gondor on liquidation mechanism, liquidation threshold, early closure, and the YES--NO merge. VANTA trades against Gondor on legibility (formula vs.\ buckets), cluster-risk handling (per-market vs.\ correlated curation), and recoverability (TEE-owned vs.\ multisig). Each trade has a concrete mitigation; each mitigation has a cost in complexity and in lender-onboarding work.

% =========================================================
\section{Limits of the Autonomy Claim}
\label{sec:limits}

Four places the autonomy framing in the abstract overreaches if read literally, and the more defensible substitutes.

\textbf{``Lenders verify the measurement instead of trusting the team.''} Cryptographically true; operationally overstated. Real lenders don't reproduce TDX builds, derive the measurement, and check the chain themselves. They trust an auditor who did. The auditor is the new curator; the curator's role has shifted from \emph{ongoing fund management} to \emph{one-time image review}. That is still a meaningful reduction --- the attack surface goes from ``team behavior over many months'' to ``one image, one audit'' --- but the trust is not eliminated. Honest framing: VANTA shrinks operational trust to one-time review; it does not eliminate trusted parties.

\textbf{``VANTA scales to every market the agent can read.''} The \S\ref{sec:onboard} envelope imposes a 14-day market-age floor, a UMA-resolver whitelist, a resolution-text-template whitelist, and a 3-markets-per-day throughput cap. Under those constraints, VANTA scales to \emph{Polymarket's standard-template subset on whitelisted resolvers}, which is a meaningful but bounded fraction of Polymarket. Honest framing: VANTA removes the curator-bandwidth bottleneck for the subset of Polymarket where the gates are well-defined. That subset grows as new templates and resolvers are timelocked-onboarded by lender governance.

\textbf{``Math, not buckets.''} The formula has three free parameters and a time constant, calibrated on a public dataset of resolved markets. That is not ``math instead of judgment'' --- it is \emph{frozen judgment from training-time, executed at run-time}. A three-bucket curator is just a model with three parameters and continuous retraining; VANTA is a continuous model with periodic timelocked retraining. The real edge is not math-vs-judgment, it is \emph{retrain-frequency $\times$ verifiability $\times$ governance friction}. Honest framing: VANTA replaces ongoing per-market judgment with a continuous pricing function whose parameters are publicly calibrated, signed-event-recorded, and timelock-governed --- reducing the rate at which discretionary decisions enter the system.

\textbf{``VANTA replaces the curator.''} On the formalizable subset it does. On the hard subset it does not, and stating otherwise overreaches. Gondor's edge on its hardest markets is not bandwidth; it is curation-as-product: years of reading UMA disputes, soft consensus on which resolvers misbehave under ambiguous events, knowing that a basket of 2028-FL-senate, 2028-OH-senate, and 2028-PA-senate is one cluster bet rather than three independent positions. The agent's three reasoning loops (\S\ref{sec:loops}) and the tag system are partial substitutes; neither is full replacement. Honest framing: VANTA owns the formalizable subset where the gates have well-defined inputs and the agent's reasoning has legible reference behavior; the hard subset stays in human hands by lender-governance upgrade, not by agent fiat. The autonomy claim is about the formalizable subset, where bandwidth is the binding constraint; on the hard subset, where curation is the product, we explicitly do not compete.

These are not retreats from the thesis. They are sharpening it. The defensible thesis is bounded verifiable autonomy on Polymarket's formalizable subset, with continuous reasoning loops above the gate floor, and attestation-liveness as the custody backstop. That is enough to differentiate from Gondor on the markets where it matters, defend in front of a serious lender, and ship.

% =========================================================
\section{Architecture}
\label{sec:arch}

Four surfaces. Three substantive (enclave, on-chain, oracle); one for verification (verifier CLI / endpoint).

\begin{figure}[h]
\centering
\begin{minipage}{0.96\textwidth}
\small\ttfamily
\begin{verbatim}
  +-- EigenCompute TDX enclave -------------------------------------------+
  |  VANTA the agent:                                                     |
  |    Ed25519 signing key (RAM, never leaves)                            |
  |    HKDF-derived wallets: treasury, origination, vault-operator        |
  |    Three loops: credit, model, operational                            |
  |    Six-hour onboarding cycle: gates -> OnboardEvent                   |
  |  HTTP surface: /api/tee  /api/attestation  /api/events/:id            |
  |                /api/origination  /api/loans/:loanId/settle            |
  |                /api/onboarding/candidates  /api/onboarding/decisions  |
  +----+----------------------------+--------+----------------------------+
       | signed events              | on-chain txs
       v                            v
  +-- Off-chain log -----------+  +-- On-chain ----------------------+
  | Walrus + IPFS (full log)   |  | Base Sepolia:                    |
  | Weekly tip anchor on Base  |  |   LpVault    0xD9AFC3...0f076    |
  | Content-addressed events   |  |   LoanBook   0xA6004c...a876f    |
  +----------------------------+  |   AnchorRegistry (weekly tip)    |
                                  |   OnboardingRegistry             |
                                  |                                  |
                                  | Polygon Amoy:                    |
                                  |   Polymarket CTF (positions)     |
                                  |   VantaVault 0xD9AFC3...0f076    |
                                  +----------------------------------+

  +-- Polymarket feed ----+      +-- Verifier --------------------------+
  | CLOB depth, TWAP      |      | vanta-verify <url> --rpc <base-rpc>  |
  | UMA dispute history   |      |   reads /api/events; walks to        |
  | Resolution criteria   |      |   genesis; cross-checks paramsHash   |
  +-----------------------+      |   against LoanBook on Base Sepolia.  |
                                 +--------------------------------------+
\end{verbatim}
\end{minipage}
\caption{Four surfaces. The enclave signs and decides; the on-chain side executes; the Polymarket feed is the input boundary; the verifier is the lender-side audit interface.}
\label{fig:arch}
\end{figure}

% =========================================================
\section{Threat Model}
\label{sec:threats}

\subsection{Defended}

\begin{center}\small
\begin{tabular}{@{}p{0.42\linewidth}p{0.52\linewidth}@{}}
\toprule
\textbf{Attack} & \textbf{Defense} \\
\midrule
Extract the signing key & Key in enclave RAM; never persisted. \\
Silent risk-model change & Every change is a signed promotion event with a 7-day timelock. \\
Fabricate a loan & Every credit event cross-references an on-chain tx; \texttt{paramsHash} bound at origination. \\
Drain the lender vault & Outflow requires a VANTA signature on a constitutional schedule; cap envelope bounds new-market exposure. \\
Forge VANTA's signature & Ed25519 unforgeable without the key. \\
Inflate pledged position & CTF balance read on-chain pre-origination. \\
LLM-provider compromise (resolution-text gate) & Three-role rotation across providers; LLM output is a gate, not an upgrade --- failed parses default-reject. \\
Onboard a garbage market & Seven-gate threshold + 3/24h throughput cap + \$25k probation envelope + UMA-resolver and template whitelist. \\
Pump-and-borrow on thin CLOB & $\min(\text{best-bid}, \text{TWAP})$ pricing; depth gate at \$250k 5\%-impact; stale-feed circuit breaker. \\
\bottomrule
\end{tabular}
\end{center}

\subsection{Not defended, disclosed explicitly}

\textbf{T1. \texttt{upgradeApp} residual.} EigenCompute derives \texttt{MNEMONIC = HMAC(appId)} against a vendor-held root~\cite{eigencompute-arch}. Whoever holds \texttt{upgradeApp} can publish a new image with the same HKDF secrets. Mitigation: fresh Ed25519 per boot makes a malicious image fork the attestation chain visibly; stewards multisig + 7-day timelock; attestation-liveness exit caps the damage at a graceful unwind.

\textbf{T2. KMS-operator compromise.} The KMS that issues the attested \texttt{MNEMONIC} is operated by the cloud provider. Disclosed; not in VANTA's control. The TVL ceiling for VANTA stays conservative until KMS is multi-vendor.

\textbf{T3. TEE.fail.} DDR5 memory-bus interposition~\cite{teefail} unpatched by Intel and AMD. Physical datacenter security is the residual defense.

\textbf{T4. Oracle failure.} UMA dispute rate $\approx 1.3$--$1.67\%$~\cite{uma-disputes}; whitelisted-proposer regime introduced August 2025 covers 96\% of Polymarket proposals at 99.7\% accuracy~\cite{uma-managed}. Priced into $\gamma\rho$; per-market concentration capped; resolver whitelist required for onboarding.

\textbf{T5. Adverse selection.} Borrowers select into mispriced loans. Adversary research role (model loop) is the first defense; gate thresholds the second; per-market cap the third (any single mispricing is bounded).

\textbf{T6. Liquidity at liquidation.} 10-min auction, 20-bps spread floor above 30-min TWAP. Underlying YES--NO merge mechanism is Polymarket-structural; we inherit Polymarket's liquidation latency, not improve it.

\textbf{T7. Cluster risk.} Per-market view misses correlated risk across markets sharing a resolution event (e.g.\ all 2028-election sub-markets). Two layers of mitigation: the tag system inflates risk for tagged clusters as a contract-enforced floor; the credit loop (\S\ref{sec:loops}) reasons about correlation patterns above the floor, including correlations the pre-registered tag set did not anticipate. Honest acknowledgment: both layers are weaker than a domain expert who has been watching a resolver for years and knows the whole cluster is one bet. The reduction VANTA actually offers is that the agent's cluster reasoning is captured in signed traces a lender can audit, where the human curator's was not. Reasoning is not stronger than the human; it is more legible.

\textbf{T8. Calibration as governance surface.} Formula parameters $(\alpha, \beta, \gamma, \tau_{\mathrm{gap}})$ are calibrated periodically on a public dataset; calibration commits go through 4-of-7 lender-quorum + 7-day timelock. The defense is not ``no humans''; it is ``humans only choose between calibrations reproducible from public data, and they cannot touch funds.''

\textbf{T9. EigenCompute deprecation.} If Eigen deprecates the TDX image, churns operators, or the attestation chain breaks, the admin key cannot be reproduced. Funds would be stranded behind a contract no one can call. Mitigation: Invariant~5 (attestation-liveness exit) converts this from ``stranded'' to ``permissionless graceful unwind.'' Cost: a 14-day Eigen outage forces a protocol pause.

% =========================================================
\section{Economics}
\label{sec:econ}

Borrower APR is 10--25\% for conservative loans, 20--35\% for long-tenor. Gross spread, after realized loss-given-default, is 8--14\%. Curator (where one exists --- Morpho-style external wrapper of VANTA's risk engine) takes 10--20\% of net; VANTA earns 0.5--1.0\% of TVL in oracle and origination fees. At $\sim$\$5M TVL with a 0.75\% blended fee, fee income covers enclave hosting (\$300--800/month for EigenCompute), LLM inference (\$100--400/month at gate-cycle frequency), gas (\$200--600/month at current Base Sepolia/Polygon Amoy levels), and oracle reads with surplus.

The autonomy framing has a direct economic implication: VANTA does not pay a curator. Gondor's three-fund model implies a team curating, rebalancing, and onboarding --- that is salaried work, recovered from the spread. VANTA recovers the same spread; the line item that was salary is now (a) image audit (one-time, contracted), (b) governance multisig (volunteer or compensated, but not full-time), and (c) enclave infrastructure (cloud-vendor cost). The economics scale linearly in TVL with no proportional headcount cost.

\textbf{Fee schedule.} Origination fee: 50 bps of principal, paid by borrower at origination. Oracle fee: 25 bps of TVL annualized, paid by lenders. Onboarding fee: \$0 to the borrower (eliminating a friction Gondor charges by curator-bandwidth); cost is internalized in the rate. Bond requirements: \$1k to flag a market (refundable on upheld flag, slashable on frivolous).

\textbf{Tranching.} Future work, not v1. A senior/junior split on the single LpVault preserves the fund-segmentation UX Gondor offers (Conservative/Moderate/Growth) without curator labor: junior shares absorb first loss; senior shares earn a lower, capped yield. We do not ship this in v1; we acknowledge it is the right answer for retail lender onboarding once v1 has proved out.

% =========================================================
\section{Roadmap}
\label{sec:roadmap}

The plan is to land each section of this paper as a shipped artifact, in order of risk-to-lender, before opening the protocol to depositors who are not the team.

\textbf{Now (shipped).} \texttt{LpVault} and \texttt{LoanBook} on Base Sepolia, \texttt{VantaVault} on Polygon Amoy --- all owned by the in-enclave admin address \texttt{0x8275...26e6}. EigenCompute App ID \texttt{0xaa4875Ff...} runs the haircut formula and origination loop. Eleven TypeScript packages (tee, events, treasury, constitution, venue-poly, mark, pledge, haircut, origination, verify, runtime). Demo runs locally on anvil forks in 90 seconds; \texttt{vanta-verify} walks the chain to genesis.

\textbf{Month 2 --- onboarding gates and image audit.} Implement the seven-gate cycle of \S\ref{sec:onboard}, with the throughput cap and probation envelope enforced contract-side. Publish the image build with reproducible-build instructions, request audit from one of \{Trail of Bits, OpenZeppelin, Code4rena\}. Auditor signs the image measurement; signature is committed in the contract. The first public depositor checks the auditor's signature, not the image themselves.

\textbf{Month 3 --- attestation-liveness exit and governance multisig.} Ship Invariant~5 to contract: 14-day silence triggers Withdraw-Only mode; permissionless withdrawal at TWAP-marked NAV. Deploy 4-of-7 lender-quorum multisig for emergency pause and parameter calibration commits. Calibration v0 published with public 12-month replay dataset and reproducible calibration code.

\textbf{Month 4 --- mainnet, with caps.} Move to Base mainnet with a hard \$1M TVL ceiling. First 30 days: lender-quorum may pause onboarding without notice; \$25k probation envelope kept; bond-and-flag system live. If no T1--T9 events trigger, ramp ceiling.

\textbf{Month 5 --- multi-position borrower UX.} Ship the batched-origination event schema; deploy a position-portfolio view served from the same signed event log the verifier reads; extend the credit loop (\S\ref{sec:loops}) with an exit planner that schedules correlated settlements in a single transaction. A holder of ten Polymarket positions originates ten loans in one user action and settles correlated positions in batches as their cluster events resolve. VantaVault tracks position-level escrow already, so this is a user-facing surface plus an agent-side scheduler over state the credit loop already maintains, not a new contract subsystem.

\textbf{Month 6+.} Tranches. Cross-margining (Gondor V1's diversified-collateral model, but with formula-derived correlation haircuts). Whitelisted Hyperliquid HIP-4 markets~\cite{hyperliquid-hip4} when their resolution criteria stabilize. Kalshi event-contract integration is gated on a separate FCM-compliance review and is not in current scope.

The ordering is not arbitrary. Each step closes a specific lender-grade question: \emph{what custody risk am I taking?} (image audit), \emph{what happens if Eigen disappears?} (attestation-liveness exit), \emph{who can pause this and why?} (lender-quorum multisig), \emph{what's my downside on a bad onboarding?} (probation envelope). A serious lender depositing in Month 4 has answers to all four.

% =========================================================
\section{Related Work}
\label{sec:related}

\textbf{Polymarket lending stack.} PolyLend~\cite{polylend} is Polymarket's own peer-to-peer lending product, dormant at \$156 TVL~\cite{polylend-tvl}. Gondor.fi~\cite{gondor-docs,gondor-seed} is the curated-team beta this paper most directly contrasts with. Both build on the Polymarket CTF proxy-wallet pattern~\cite{polymarket-proxy}. Morpho~\cite{morpho} provides the underlying lending primitive Gondor uses; we do not use Morpho for v1 (the LpVault/LoanBook design is bespoke to keep custody enclave-bound).

\textbf{Sovereign-agent infrastructure.} EigenCloud AgentKit~\cite{eigenagentkit} for the TEE-attested, on-chain-paying agent platform. Coinbase AgentKit~\cite{coinbaseagentkit} for the wallet+action layer. bobIsAlive~\cite{bobisalive} (Synthesis Hackathon 2026 winner) and Sovra~\cite{sovra} as prior projects from this author shipping enclave-derived custody and self-paying business primitives; VANTA reuses the substrate, not the use case.

\textbf{Prediction-market credit math.} Kirilenko et al.~\cite{kirilenko2025} for $\sigma(\tau)$ under jump-diffusion. Wolfers and Zitzewitz~\cite{wolfers-zitzewitz} for the martingale baseline. Vaugirard~\cite{vaugirard2003} and Chao~\cite{chao2018} for catastrophe-bond pricing relevant to the gap term. L\'opez de Prado~\cite{lopezprado} for backtest discipline behind the model loop.

\textbf{Adjacent prediction-market venues.} Hyperliquid HIP-4~\cite{hyperliquid-hip4} (testnet March 2026, no mainnet date as of writing) introduces outcome contracts and may become a second venue for VANTA after resolution criteria stabilize. Kalshi's affiliate Kinetic Markets received NFA FCM approval March 2026~\cite{kalshi-fcm}, clearing margin trading on event contracts under CFTC rule; that is parallel infrastructure under different regulatory regime, not on-chain credit.

\textbf{Oracle layer.} UMA Optimistic Oracle~\cite{uma-disputes,uma-managed} is the resolution rail for Polymarket; the whitelisted-proposer regime introduced August 2025 (37 proposers, 99.7\% accuracy on 96\% of proposals) materially reduces the dispute-rate input to the haircut formula's $\gamma\rho$ term.

% =========================================================
\section{Conclusion}
\label{sec:conclusion}

VANTA is two things at once.

It is an honest narrowing of Gondor: autonomous on the subset of Polymarket markets where the gates are well-defined, manual everywhere else by deferring those markets to a future timelocked governance upgrade. The narrowing is a feature. It is what makes the autonomy claim defensible.

It is a verifiability claim about the lending stack: the engine's code is an attested image; parameters and the agent's three reasoning loops emit signed traces lenders can audit; the contracts are owned by a key only that image can derive; a 14-day silence triggers withdraw-only without admin involvement. The team's job becomes publishing an auditable image, not managing a fund. Lender-grade trust shrinks from operational to one-time review.

What remains: ship the onboarding gates (Month 2), open the auditor pass (Month 2), wire the attestation-liveness exit (Month 3), deploy the lender-quorum multisig (Month 3), move to mainnet under a TVL ceiling (Month 4). Each step closes a specific lender-grade question. None of them require a curator.

We invite critique --- on the formula, on the gates, on the threat model, on the limits in \S\ref{sec:limits}. The thesis is sharp because we have stated it sharply, and stated also where it does not hold. The work from here is making the bands fill.

% =========================================================
\begin{thebibliography}{99}
\footnotesize

\bibitem{polymarket-volume} BitKE. \emph{Polymarket crosses \$10B monthly volume in March 2026.} April 2026. \url{https://bitcoinke.io/2026/04/polymarket-in-march-2026/}.

\bibitem{polymarket-oi} CryptoTimes. \emph{Polymarket TVL crosses \$500M following CLOB v2 rollout.} April 2026.

\bibitem{polylend} Polymarket. \emph{PolyLend: peer-to-peer lending against CTF positions.} Source repository, Polymarket organization on GitHub.

\bibitem{polylend-tvl} DefiLlama. \emph{PolyLend TVL.} Accessed April 2026.

\bibitem{gondor-seed} Gondor.fi. \emph{Announcement of \$2.5M pre-seed round with Prelude Ventures, Castle Island Ventures, and Maven11.} December 2025.

\bibitem{gondor-docs} Gondor.fi. \emph{Protocol documentation.} 2025--2026.

\bibitem{morpho} Morpho Labs. \emph{Morpho lending protocol.}

\bibitem{theowhale} Chainalysis. \emph{Reconstructing the Th\'eo wallet cluster: \$85M across 11 accounts on Polymarket's 2024 election market.} 2025.

\bibitem{predictionmarkets-1b} Aggregated open interest across Polymarket, Kalshi, and on-chain prediction venues, April 2026.

\bibitem{kalshi-fcm} National Futures Association. \emph{Kinetic Markets, LLC --- Futures Commission Merchant approval.} March 2026.

\bibitem{polymarket-proxy} Polymarket. \emph{Gnosis Safe proxy wallet integration for CTF positions.} Polymarket developer documentation.

\bibitem{kirilenko2025} A.~Kirilenko et al. \emph{Jump-diffusion calibration for binary event contracts.} Working paper, 2025.

\bibitem{uma-disputes} UMA Protocol. \emph{Optimistic Oracle dispute statistics, 2024--2025.}

\bibitem{uma-managed} UMA Protocol. \emph{Whitelisted-proposer regime: 96\% of Polymarket proposals at 99.7\% accuracy.} August 2025.

\bibitem{eigencompute-arch} EigenLayer. \emph{EigenCompute architecture and KMS design.} EigenCloud documentation, 2026.

\bibitem{teefail} TEE.fail collaboration. \emph{DDR5 memory-bus interposition attacks against Intel TDX and AMD SEV-SNP.} 2024--2025.

\bibitem{hyperliquid-hip4} Hyperliquid. \emph{HIP-4: Outcome contracts on Hyperliquid.} Testnet, March 2026.

\bibitem{eigenagentkit} EigenLayer. \emph{EigenCloud AgentKit.} 2025--2026.

\bibitem{coinbaseagentkit} Coinbase. \emph{AgentKit: wallet and action layer for autonomous agents.} 2025.

\bibitem{bobisalive} W.~Okechukwu. \emph{bobIsAlive: a sovereign agent on EigenCompute.} Synthesis Hackathon 2026 winner.

\bibitem{sovra} W.~Okechukwu. \emph{Sovra: enclave-derived custody and self-paying business primitives.} 2026.

\bibitem{wolfers-zitzewitz} J.~Wolfers and E.~Zitzewitz. \emph{Prediction markets.} Journal of Economic Perspectives, 18(2):107--126, 2004.

\bibitem{vaugirard2003} V.~Vaugirard. \emph{Pricing catastrophe bonds by an arbitrage approach.} Quarterly Review of Economics and Finance, 43(1):119--132, 2003.

\bibitem{chao2018} W.~Chao. \emph{Valuation of catastrophe equity put options with stochastic interest rates.} Journal of Risk and Financial Management, 11(2):26, 2018.

\bibitem{lopezprado} M.~L\'opez de Prado. \emph{Advances in Financial Machine Learning.} Wiley, 2018.

\end{thebibliography}

\end{document}
